\bookchapter{נספח F: נושא רכישה ושימוש במסך מתקפל}

\chapterintro{נתחיל במסקנה:}

\booksection{לאיזה שלב הגיעו הטלפונים המתקפלים?}

. הם נמצאים בין מוצר נישה לחובבים למוצר צריכה המוני, ומתאימים לקהל מסוים.\textbf{כיום יש שתי קטגוריות עיקריות:}קיפול אופקי (סגנון ספר)

: כשהוא סגור הוא בגודל טלפון רגיל, כשהוא פתוח הוא טאבלט קטן כמעט מרובע. דוגמאות: סדרת Huawei Mate X, סדרת Samsung Galaxy Z Fold, סדרת Xiaomi MIX Fold, סדרת vivo X Fold. היתרון העיקרי: “טלפון אחד = טלפון + טאבלט”.

\begin{itemize}
  \item \textbf{קיפול אנכי (סגנון צדפה)}: כשהוא סגור הוא נראה כמו קופסת פודרה, כשהוא פתוח הוא טלפון בגודל רגיל. דוגמאות: סדרת Huawei Pocket, סדרת Samsung Galaxy Z Flip, OPPO Find N Flip, Xiaomi MIX Flip. היתרון העיקרי: "קומפקטי ונייד + הצצה מהירה במסך החיצוני".
  \item \textbf{התקדמויות משמעותיות ב-2024-2025: חיי הציר מדורגים בדרך כלל ל-500,000 קיפולים (100 פתיחות וסגירות ביום יספיקו ל-13 שנים), זכוכית UTG דקה במיוחד מחליפה את ציפוי הפלסטיק המוקדם, קמט הקיפול עדיין נראה לעין אך כמעט לא מורגש במגע, והמחיר ירד מרמת 10,000 יואן לרמת 5,000-6,000 יואן.}שיפוט במשפט אחד
\end{itemize}

2024-2026: ציר 500,000, UTG, 4000-5000.

\begin{tcbinfo}{2. למי מתאים לקנות טלפון מתקפל?}
על סמך משוב משתמשים אמיתי, הסבירות שאנשים מהסוגים הבאים יאהבו טלפון מתקפל גבוהה:

\end{tcbinfo}

\booksection{למי מתאים לקנות טלפון מתקפל?}

: אנשים שמרבים לקרוא PDF, להעיר על מסמכים ולקרוא טקסטים ארוכים בטלפון. המסך הפנימי בגודל קרוב ל-7.5 אינץ׳ במצב פתוח מאפשר קריאת PDF בגודל A4 ללא צורך בהגדלה — שינוי איכותי בחוויה.

\begin{enumerate}
  \item \textbf{משתמשים עסקיים שמרבים בנסיעות/פגישות}: לא רוצים לסחוב טאבלט אבל מרגישים שמסך הטלפון קטן מדי. טלפון מתקפל אופקי מאפשר צפייה במצגות PPT, חתימה על חוזים והשוואת מיילים בשני מסכים — יעילות גבוהה בהרבה מטלפון רגיל.
  \item \textbf{אנשים שאוהבים ליצור במצב "חצובה"}: קיפול ל-90 מעלות על השולחן — חצי עליון למסגור וחצי תחתון ללחיצה על הצמצם (סלפי/צילום זמן-שגב), חצי עליון לווידאו וחצי תחתון לרשימות, חצי עליון לשידור חי וחצי תחתון לקריאת תגובות — זה משהו שטלפון רגיל לא יכול לעשות.
  \item \textbf{משתמשות שרוצות גם ניידות וגם מסך גדול}: טלפון מתקפל אנכי סגור נכנס לתיקון קטן, פתוח הוא טלפון רגיל, והמסך החיצוני יכול לשמש גם כמראת איפור ולהצגת התראות.
  \item \textbf{האנשים האלה לא צריכים לקנות}מי שמשתמש בטלפון רק לסרטונים קצרים, וואטסאפ ושיחות טלפון — טלפון מתקפל לא נותן לך שום יתרון, להפך — הוא כבד ועבה יותר.
\end{enumerate}

\begin{tcbwarn}{מי שיש לו ידיים קטנות או דלקת גידים בפרק כף היד — טלפון מתקפל אופקי שוקל בדרך כלל 230 גרם ויותר, קשה להחזיק ביד אחת.}
\begin{itemize}
  \item מי שתקציבו מוגבל וצריך טלפון ל-4 שנים ויותר — ערך המכירה החוזרת של טלפון מתקפל נמוך משל טלפון רגיל, ועלות התיקון גבוהה.
  \item 3. נקודות מפתח לבחירה: 6 מדדים שצריך לבדוק
  \item אל תתבלבל מהמילה "מתקפל", בעת הבחירה התמקד ב:
\end{itemize}

\end{tcbwarn}

\booksection{נקודות מפתח לבחירה: 6 מדדים שצריך לבדוק}

: ציר טיפת מים (קמט רדוד, חיים ארוכים, הנפוץ כיום) עדיף על ציר בצורת U (קמט עמוק, טכנולוגיה מוקדמת). שאל את המוכר או בדוק בסקירות באיזה סוג מדובר.

\begin{itemize}
  \item \textbf{חומר המסך הפנימי}: חייב להיות זכוכית UTG דקה במיוחד, לא ציפוי הפלסטיק המוקדם. UTG מרגיש כמו זכוכית, בעוד שציפוי פלסטיק יוצר שקע כשמעבירים עליו ציפורן ונוטה להישרט.
  \item \textbf{מראה הקמט}: לך לחנות פיזית ופתח במו ידיך. הרמה המקובלת ב-2025: כמעט לא נראה מלפנים כשהמסך דולק, באור אלכסוני רואים חריץ רדוד, באצבע מרגישים שקע קל אבל לא חד.
  \item \textbf{גודל המסך החיצוני}:בחנות. 2026: כמעט לא רואים.
  \item \textbf{משקל ועובי}: בטלפון מתקפל אופקי 230-260 גרם זה הנורמה, פחות מ-230 גרם נחשב מצוין; עובי במצב סגור 12-14 מ"מ, פחות מ-12 מ"מ נחשב דק. בטלפון מתקפל אנכי, הוא צריך להיכנס בקלות לכיס חולצה כשהוא סגור.
  \item \textbf{שירות ואחריות}: התמקד במחיר תיקון מסך פנימי שבור. חלק מהמותגים מציעים החלפת מסך פנימי חינם בשנה הראשונה, אחרים גובים 2,000-4,000 יואן. זוהי העלות הסמויה הגדולה ביותר של טלפון מתקפל.
  \item \textbf{רשימת בדיקה בחנות פיזית}פתח וסגור 10 פעמים במו ידיך, הרגש אם הציר נע חלק והאם יש רעשים חריגים.
\end{itemize}

\begin{tcbpractice}{לאחר הפתיחה, הסתכל על הקמט מול אור אלכסוני, ואז מלפנים כשהמסך דולק.}
\begin{enumerate}
  \item החזק ביד אחת את הטלפון הפתוח למשך 30 שניות, הרגש אם המשקל נסבל.
  \item נסה מצב חצובה: קפל ל-90 מעלות על השולחן, נסה לצלם ולצפות בווידאו.
  \item ברר היטב את מחיר תיקון המסך הפנימי ומדיניות האחריות.
  \item 4. טיפים ייחודיים לשימוש בטלפון מתקפל
  \item אחרי הקנייה, אלו הדברים שלא תוכל לחוות בטלפון רגיל:
\end{enumerate}

\end{tcbpractice}

\booksection{טיפים ייחודיים לשימוש בטלפון מתקפל}

בטלפון מתקפל אופקי פתוח, רוב המערכות תומכות בפיצול מסך למעלה/למטה או שמאל/ימין. תרחישים אופייניים: חצי עליון לשיעור וידאו, חצי תחתון לרשימות; וואטסאפ משמאל ומסמך מימין; מפת ניווט למעלה ומוזיקה למטה. הרבה יותר שימושי מפיצול מסך בטלפון רגיל כי המסך גדול מספיק וכל חצי לא צפוף.

\booksubsection{מצב חצובה (קיפול 90 מעלות)}

זהו שימוש ה“קלף מנצח” של הטלפון המתקפל:

\booksubsection{סלפי/תמונה קבוצתית}

: הנח על השולחן בקיפול 90 מעלות, חצי עליון למסגור, חצי תחתון ללחיצה על הצמצם, אין צורך לבקש מאף אחד לצלם.

\begin{itemize}
  \item \textbf{צילום זמן-שגב/צילום לילה}: באותה צורה, צילום יציב על השולחן, חוסך חצובה.
  \item \textbf{שיחות וידאו}: ידיים חופשיות, והזווית שבה רואים אותך טבעית יותר.
  \item \textbf{צפייה בווידאו}: על השולחן כמו טלוויזיה קטנה, שימושי במיוחד על שולחן מתקפל במטוס/רכבת מהירה.
  \item \textbf{תצוגה מקדימה במסך חיצוני + צילום במסך פנימי}: כשמצלמים את עצמך במצלמה האחורית, המסך החיצוני מראה תמונה בזמן אמת, הקומפוזיציה מדויקת יותר.
  \item \textbf{3. הצצה מהירה במסך החיצוני}בטלפון מתקפל אנכי סגור, המסך החיצוני מאפשר: צפייה בשעה/התראות/מזג אוויר, מענה בוואטסאפ (הקלטה לתמלול או תשובה מהירה), כרטיס תחבורה ציבורית, ניווט, מראת איפור (באמצעות המצלמה הקדמית). אין צורך לפתוח בכל פעם, חוסך בסוללה ובטרחה.
\end{itemize}

\booksubsection{מצב מחשב}

חלק מהטלפונים המתקפלים האופקיים (Huawei, Samsung) מאפשרים לעבור ל“מצב מחשב” במצב פתוח, עם חיבור למסך חיצוני, מקלדת ועכבר זה עובד כמו מחשב פשוט. מספיק לעיבוד מסמכים זמני ומענה למיילים. לא כלי פרודוקטיביות עיקרי, אבל שימושי מאוד לשעת חירום.

\booksubsection{נקודות תחזוקה: איך לגרום לטלפון המתקפל להחזיק יותר זמן}

הדבר העדין ביותר בטלפון מתקפל הוא המסך הפנימי והציר. ההרגלים הבאים יכולים להאריך משמעותית את חייו:

\booksection{נקודות תחזוקה: איך לגרום לטלפון המתקפל להחזיק יותר זמן}

: מתחת ל-0 מעלות זכוכית UTG הופכת שבירה, חמם אותה בכיס כמה דקות לפני הפתיחה. זו נקודה חשובה שכתובה במדריך הרשמי אבל אף אחד לא קורא.

\begin{itemize}
  \item \textbf{הימנע ממגע ציפורניים/חפצים קשים עם המסך הפנימי}: למרות שלמסך הפנימי יש שכבת הגנה, הוא רך יותר מהמסך החיצוני. השתמש בכריות האצבעות, אל תדקור בציפורניים.
  \item \textbf{וודא שאין חפצים זרים על המסך לפני הסגירה}: גרגר חול אחד שנתפס בסגירה עלול למעוך את המסך הפנימי. העיף מבט לפני הסגירה.
  \item \textbf{אל תדביק מגן זכוכית רגיל}: המסך הפנימי תומך רק במגן הרשמי או מגן UTG ייעודי. מגן זכוכית רגיל יפגע בקיפול ובמגע. מומלץ לגשת לשירות הרשמי להדבקה.
  \item \textbf{מנע כניסת אבק ונוזלים לציר}: למרות שרוב הדגמים מדורגים IPX8 לעמידות במים, דירוג עמידות האבק בדרך כלל נמוך. היזהר בחוף ים, אתרי בנייה וסביבות דומות.
  \item \textbf{פתח וסגור בכוח אחיד}: אל תפתח בטריקה ביד אחת, פתח בשתי ידיים מהאמצע בתנועה אחידה — זו הדרך הבטוחה ביותר.
  \item \textbf{תרחישים נפוצים לנזק למסך פנימי}סגירה על כבל אוזניות/כבל טעינה — ימעוך את המסך הפנימי ישירות.
\end{itemize}

\textbf{בתיק עם מפתחות — המסך החיצוני בסדר, אבל המסך הפנימי נסדק מהמפתחות.}\par

\begin{itemize}
  \item בחורף, מעבר מקור לחדר חם ופתיחה מיידית — הפרש טמפרטורות גורם לשינוי במתח הזכוכית.
  \item אחרי שנה של שימוש, תחושה ש"הקמט העמיק" וניסיון להסיר את המגן המקורי — הסרת המגן תפגע גם במסך הפנימי.
  \item 6. ה"מהמורות" של טלפונים מתקפלים וציפיות ריאליות
  \item לפני הקנייה, קח בחשבון את הדברים הבאים:
\end{itemize}

\booksection{ה“מהמורות” של טלפונים מתקפלים וציפיות ריאליות}

: הטכנולוגיה של 2025 לא יכולה להשיג “העלמה מוחלטת”, רק “כמעט לא נראה מלפנים כשהמסך דולק”. אם אתה פרפקציוניסט וכל יום תתעצבן מהקמט — אל תקנה.

\begin{itemize}
  \item \textbf{כבד ועבה יותר}:לא "בלי עקבות". פרפקציוניסט? לא.
  \item \textbf{חיי סוללה בדרך כלל בינוניים}: מנגנון הציר תופס מקום, מרחב הסוללה מצומצם, חיי הסוללה בדרך כלל פחות טובים מטלפון רגיל באותו מחיר.
  \item \textbf{ערך מכירה חוזרת נמוך}: אחרי שנה, מחיר יד שנייה הוא בדרך כלל 50-60\% ממחיר חדש, יורד מהר יותר מטלפון רגיל.
  \item \textbf{תיקון יקר}: תיקון מכלול מסך פנימי עולה 2,000-4,000 יואן, תיקון ציר 500-1,500 יואן. ברר היטב את מדיניות האחריות לפני הקנייה.
  \item \textbf{התאמת אפליקציות לא אחידה}: אפליקציות פופולריות כבר מותאמות למסך גדול, אבל אפליקציות נישה עלולות פשוט "להימתח", חוויה בינונית.
  \item \textbf{טלפון מתקפל הוא לא "טלפון טוב יותר", אלא "טלפון בצורה אחרת". קנה אותו כי אתה צריך את הצורה הייחודית שלו, לא כי הוא "מתקדם יותר".}—— על הפילוסופיה של בחירת טלפון מתקפל
\end{itemize}

7. איך לבחור בין המותגים המובילים (2025)

בלי דגמים ספציפיים (שמתיישנים), רק אופי המותג, כדי לעזור לך להתמקד במהירות:

\booksection{ז', מותגים (2026)}

: אמת המידה לקיפול אופקי, אינטגרציה חזקה עם מערכת HarmonyOS, יחס מסך חיצוני קרוב לטלפון רגיל, ניתן להשתמש כטלפון עיקרי כשהוא סגור. מחיר גבוה אבל ערך המכירה החוזרת הטוב ביותר יחסית.

\begin{itemize}
  \item \textbf{סדרת Samsung Galaxy Z Fold/Flip}: חלוצת הטלפונים המתקפלים, כיוון ציר בשל, Z Flip הוא המוביל בעיצוב בקיפול אנכי. חיסרון — מעט מרכזי שירות בסין.
  \item \textbf{סדרת Xiaomi MIX Fold/Flip}: גישת תמורה למחיר, מפרט עשיר, מחיר יחסית נגיש. מתאים למשתמשים שרוצים לנסות בלי להוציא יותר מדי.
  \item \textbf{סדרת vivo X Fold}: יחס המסך החיצוני בקיפול אופקי הקרוב ביותר לטלפון רגיל, חוויית שימוש טובה כטלפון רגיל כשהוא סגור.
  \item \textbf{סדרת OPPO Find N}: מידות קומפקטיות יחסית בקיפול אופקי, ניתן לאחיזה ביד אחת, מתאים למשתמשים שחושבים שקיפול גדול כבד מדי.
  \item \textbf{סדרת Huawei Pocket}: המסך החיצוני הגדול והשימושי ביותר בקיפול אנכי, פונקציות ידידותיות לנשים כמו מראת איפור ובדיקת הגנה מקרינת שמש עשויות בקפידה.
  \item \textbf{המלצת בחירה: תקציב גבוה ומערכת אקולוגית — Huawei Mate X; הציר הבשל ביותר — Samsung Z Fold; תמורה למחיר — Xiaomi/vivo; קיפול אנכי — Huawei Pocket ו-Samsung Z Flip.}תזכורת אחרונה
\end{itemize}

טכנולוגיית הטלפונים המתקפלים מתעדכנת מדי שנה, נספח זה נכתב ב-2025. לפני הקנייה בדוק תמיד סקירות עדכניות ומשוב משתמשים, עם דגש על: חיי ציר בפועל, מראה הקמט, ומדיניות תיקון מסך פנימי. לגעת במו ידיך בחנות תמיד חשוב יותר מלקרוא סקירות.

\begin{tcbinfo}{תזכורת אחרונה}
2026. בדוק: ציר/קמט/תיקון. נסה.

\end{tcbinfo}
